\chapter{Objections and Replies}\label{app:objections}

The strongest objections deserve the strongest form of themselves before any reply. What follows states five of them at full strength, then answers each from the same three things the rest of the book rested on: the bedrock, the postulate, and finitude. None of the replies is propped up from outside. If an objection cannot be met from the structure, it should stand.

\emph{This is anti-realism, or instrumentalism, in disguise.} The book keeps saying that no cut is privileged, that realism is a cut, that even the self is a cut. Strip the vocabulary and this is the old idealist move: nothing is really out there; it is all observer-relative construction.

The postulate denies none of that reality; it declines to privilege one decomposition of it as the correct one. The book is emphatically realist about the structure. The totality of computable possibility, the measure on it, and the theorems that hold of it are not observer-relative. What the book refuses is the further claim that one carving, into ``the one actual world,'' into systems, into selves, is the structure's own joint rather than a boundary an observer drew. That is not ``nothing is real.'' It is ``the structure is real and has no privileged parts.'' And realism about a single actual world is not the neutral default it feels like; it is itself a cut, and specifying which world costs information. The book holds the existence question open, undecidable from inside, and asserts neither that the carved things are real nor that they are not. Anti-realism asserts they are not. The book asserts neither, which is the more disciplined place to stand.

\emph{The universal prior is uncomputable, so it is idle.} $M$ cannot be computed; $K(x)$ cannot be computed. A quantity no agent can ever calculate is a curiosity, not a foundation, and a worldview built on it is built on sand.

Uncomputability does not make $M$ idle. It makes it a standard. The dominance and invariance theorems (Appendix~\ref{app:bedrock}) are statements about the ideal object, and they are exactly what give computable approximations their warrant: whatever predictor you can actually build, its guarantee is stated relative to $M$. A reference mass you cannot manufacture perfectly is not idle; it is what every real instrument approximates and is measured against. And finitude explains why no observer needs to compute the prior at all. A finite cutter is forced toward cheap cuts, so it approximates the universal prior with short, local descriptions whether it prefers to or not. The uncomputable object sets the target; finitude explains the approximation; the distance between them is not a hole in the argument but its subject. The full development is in Hutter and in Li and Vit\'anyi.

\emph{Simplicity is smuggled in as a value.} The book leans on simplicity everywhere: shorter is more probable, the simple system is the one you find yourself in. But ``prefer the simple'' is an aesthetic taste, quietly installed as though it were mathematics.

Simplicity here is not a value the book imposes; it is a scarcity the book reports. Kraft's inequality (Appendix~\ref{app:bedrock}) says the weights of a prefix-free code sum to at most one, so short descriptions are necessarily few: there is not enough probability mass for them to be common. A finite observer, with finite mass to allocate, is therefore compelled to spend most of it on short descriptions. That is finitude, not preference. The book is careful never to say the structure prefers the simple, because it has no joints to prefer with; it says only that a finite cutter cannot avoid cheap cuts. Remove the author's taste entirely and the bias remains, because it was never taste. It was counting.

\emph{The turn to care is a non-sequitur from the math.} Fifteen chapters of cold structure, and then, with no valid inference, the book turns to kindness and closes on ``Stay.'' Nothing in Bayes, Kraft, or Solomonoff entails that you should care about anyone. The ending is a sentimental addition, the very consolation the book claims to refuse.

This objection is correct, and the book says so in its own text. Care is not derived from the structure. The mathematics is silent on what to do; the structure does not require care, does not reward it, does not notice it. The turn is not an inference. It is a choice, a cut made without the structure's permission, because no permission was ever coming. The book never pretends otherwise. It calls the ending not an argument, and it asks outright, ``Is it enough? I do not know.'' The honesty is the point, not a defect to be patched. A book that claimed the mathematics entailed kindness would be consoling you with something untrue, and that is the one move this book refuses. Care is held anyway, with no floor under it, and that there is no floor is stated plainly rather than hidden. The objection names the gap. The book was already standing in it, on purpose.

\emph{The multiverse is unfalsifiable, so this is not science.} The book posits every computable world, hellscapes and Boltzmann observers among them. None of it can be tested. By the usual criterion that is not science but metaphysics in a lab coat.

The book agrees it is not a physical hypothesis, and never claimed to be. It does not assert that the multiverse exists. The existence question is held open, undecidable from inside; the book declines to claim the carvings are real. So there is no empirical bet here awaiting a telescope. What the book offers is a structural observation: that description has no privileged decomposition, that a finite observer must cut, and that ``the one real world'' is one such cut with a price. That is a claim about descriptions and the observers who make them, to be judged on whether it holds of description, not on whether an experiment could falsify it. Falsifiability is the right test for a claim about the physical world. This is not one, and it does not ask for the credit of being one.
